% ===========================================================================
%  Alden Grover — master CV (LaTeX source of truth)
% ===========================================================================
%
%  Populated from the master CV. This file is the source for BOTH published
%  artifacts:
%
%    assets/pdf/Alden_CV.pdf   the PDF (ATS / print / fidelity view)
%    _includes/cv-live.html    the in-page HTML view
%
%  Rebuild with ./cv/build.sh (see ../BUILD.md). Never hand-edit either artifact.
%
%  CONTENT RULES:
%    * This file is the source of truth. Keep it in step with _data/*.yml and
%      the _projects/ + _research/ collections, which feed the interactive tab
%      of /cv/ — the two views must not disagree.
%    * Do not invent credentials, dates, or metrics. Everything here came from
%      the master CV.
%
%  WRITING RULES (keep this file "pandoc-friendly" — see ../BUILD.md):
%    * plain `article` + geometry + hyperref (all in texlive-latex-base/-recommended).
%      Never switch to moderncv or an exotic template: pandoc's LaTeX -> HTML5
%      conversion degrades badly on those, and this file renders as HTML on /cv/.
%    * semantic \section*{} and \itemize only.
%    * no floats, minipages or multi-column tricks for anything that must
%      appear in the on-page HTML view; two-column output would also risk
%      interleaving the text layer, which ATS parsers read.
%    * dates and trailing metadata go through \cvwhen{}, never bare after
%      \hfill — pandoc silently drops \hfill, losing the separation.
%    * date ranges use a plain ASCII "-" with spaces around it, NOT "--".
%      "--" typesets an en dash (U+2013), which puts a non-ASCII character in
%      the text layer for no benefit; a literal en/em dash is worse. ASCII
%      extracts cleanly for ATS matching and avoids PDFDocEncoding surprises
%      in metadata.
% ===========================================================================

\documentclass[10pt,letterpaper]{article}

\usepackage[T1]{fontenc}
\usepackage[utf8]{inputenc}

% Margins and leading are measured against the master CV, not guessed. The
% master uses ~0.5in margins and ~13.2pt leading, and both are matched here, so
% this PDF reads at the same density as the master CV. That density no longer
% packs into 3 pages once entries are kept whole across page breaks — see the
% entry keep-together note below and the page-count note in BUILD.md §2.
%
% Rendering this content with the default 1in margins needed 5 pages — a real
% readability regression, since the master CV is the format recruiters already
% expect.
%
% History, so this is not re-tightened by mistake. An earlier revision used
% margin=0.25in + linespread=0.86: 2 pages, but at a 10.3pt pitch (vs the
% master's 13.2pt) and with under 2 lines of slack left on either page. That
% version had no room to grow — adding a single line spilled content onto a
% near-empty page 3. The values below are the LEAST aggressive that still hold
% 3 filled pages, so there is headroom again.
%
% Note this is deliberately NOT solved by shrinking the font. Dropping to a
% smaller base size would change every EC font size used, invalidating the
% fonts baked into the toolchain image (cv/warmup.tex) and costing ~1.5s of
% mktexpk on the first run after each change. Widening the text block instead
% keeps the line height and the baked fonts untouched.
\usepackage[margin=0.5in]{geometry}
% Charter (Bitstream Charter, PSNFSS bch), not Times. Charter's x-height is
% 0.482em against Times' 0.450em, so at this 10pt base it reads ~7% larger
% without a single extra page — the cheapest legibility gain available, since
% 10pt Times is the smallest-looking type in the toolchain image. Charter was
% drawn to hold up at low resolution, which is also the case here (screen and
% print at 10pt).
%
% The `helvet` line is NOT decorative: charter.sty resets \sfdefault away from
% the Helvetica that times.sty used to supply, and the T1 Computer Modern sans
% then has no scalable face in this image — headings silently become a Type 3
% font with no Unicode map (pdffonts shows "F43", uni=no), which is exactly what
% the text-layer rules below exist to prevent. Loading Helvetica keeps every
% font Type 1 and Unicode-mapped. If you swap the serif package again, check
% `pdffonts assets/pdf/Alden_CV.pdf` for a Type 3 entry before committing.
\usepackage{charter}
\usepackage{helvet}
\usepackage[hidelinks]{hyperref}

% --- ATS / copy-paste text layer --------------------------------------------
% Résumé parsers and copy-paste read the PDF text layer, and T1 (EC/LModern)
% encoding breaks it: fi and fl are emitted as single ligature glyphs with no
% text equivalent, so "file" extracts as "le", "profile" as "prole" and
% "efficient" as "ecient". That silently fails the literal matching an ATS does.
%
% Two pieces here, and which one matters is not obvious:
%
%   1. microtype's \DisableLigatures — stops pdfTeX emitting fi/fl as ligature
%      glyphs at all. This is the one that actually fixes extraction.
%   2. \pdfgentounicode — emits a ToUnicode CMap for the remaining glyphs
%      (quotes, dashes, bullets). Good practice, but not the ligature fix.
%
% Measured on this file, because the obvious approaches fail:
%   * T1 + glyphtounicode + \pdfgentounicode, no microtype -> still "prole le".
%   * Drop T1 (default OT1) -> ligatures extract fine, but accents break
%     ("Université naïve" extracts as "Universite ́ naı̈ve").
%   * T1 + \DisableLigatures -> correct accents AND clean extraction. Chosen.
%
% Do not remove this block. After editing, verify with:
%   pdftotext cv.pdf - | grep -i "prole le"    # must print nothing
%
% microtype is in texlive-latex-recommended, so this adds no new dependency,
% and it must be loaded after fontenc. \pdfgentounicode is pdfTeX-only, hence
% the guard.
%
% `expansion=false` is deliberate: microtype's font expansion needs scalable
% fonts, and on a minimal TeX install the T1 fonts are bitmap-only, so pdfTeX
% aborts with "auto expansion is only possible with scalable fonts" and
% produces NO PDF. Expansion buys subtle letter-spacing refinements; the
% build working everywhere matters more. Character protrusion (microtype's
% main optical win) stays enabled. Turning expansion back on would require a
% scalable font package such as lmodern/cmsuper.
\usepackage[expansion=false]{microtype}
\DisableLigatures{encoding = T1, family = *}
\ifdefined\pdfgentounicode
  \input{glyphtounicode}%
  \pdfgentounicode=1
\fi

% This file deliberately sticks to packages in texlive-latex-base +
% texlive-latex-recommended (fontenc, inputenc, geometry, hyperref, microtype)
% so it compiles on a minimal install — no texlive-latex-extra, and no
% texlive-fonts-extra (lmodern/cmsuper) required. Spacing is therefore done
% with plain lengths instead of the enumitem/parskip packages, which would
% otherwise pull in ~1.4 GB of dependencies for two lines of setup.

% No paragraph indent. parskip is small on purpose: with ~20 entries each
% separated by \vspace, a large parskip compounds and costs pages.
\setlength{\parindent}{0pt}
\setlength{\parskip}{0pt}
\emergencystretch=2em
% Never hyphenate: a hyphenated line-break word splits in the extracted text
% layer, and an ATS doing literal token matching misses the fragment.
\hyphenpenalty=10000
\exhyphenpenalty=10000

% Leading matched to the master CV. The 10pt base size has a 12pt default
% baselineskip, so 1.10 gives 13.2pt — the master CV measures 13.1pt at the
% same base size, i.e. the two agree. (This file was once 0.86 = 10.3pt, which
% looked cramped next to the master.)
%
% \linespread scales leading only, so no new font sizes are introduced and the
% toolchain image's baked fonts stay valid (see the geometry note above).
\linespread{1.10}

% --- PDF document metadata (ATS / recruiter tooling) -------------------------
% Application tracking systems and recruiter tools read the PDF's document
% information dictionary to answer "whose CV is this?" and "what is it?".
% Without \hypersetup, pdfTeX leaves Title/Author/Subject/Keywords EMPTY, which
% is what this file shipped before: `pdfinfo` reported blank fields on every
% build, so an automated reader had nothing to index by name.
%
% \hypersetup also removes hyperref's default "LaTeX with hyperref" authorship
% signal. Note the pdftex option is required for pdftitle/pdfauthor to take
% effect in pdfTeX; the values are plain ASCII deliberately, since PDFDocEncoding
% handles the em dash poorly in metadata strings.
%
\hypersetup{
  pdftex,
  pdftitle  = {Alden Grover - Curriculum Vitae},
  pdfauthor = {Alden Grover},
  pdfsubject= {Curriculum Vitae - Electrical and Computer Engineering},
  pdfkeywords={RTL design, ASIC design, verification, computer architecture, embedded systems, computer engineering, SystemVerilog, tape out, FPGA},
}

% Compact bullet lists, applied per-list (plain LaTeX: valid inside \itemize,
% and pandoc's LaTeX reader still sees ordinary itemize blocks).
\newcommand{\tightlist}{%
  \setlength{\itemsep}{0pt}%
  \setlength{\leftmargin}{1.0em}%
  \setlength{\parskip}{0pt}%
  \setlength{\topsep}{1.5pt}%
  \setlength{\parsep}{0pt}%
}

% --- Keeping entries whole across page breaks -------------------------------
% No entry may straddle a page break. The toolchain image carries only
% texlive-latex-base + -recommended, so \needspace is unavailable; this is a
% package-free stand-in. Called before each entry, it breaks first when fewer
% than #1 of vertical space remain, so any entry shorter than #1 is never split.
%
% #1 must be ≥ the tallest entry or a taller one can still straddle a break; the
% longest entry here (Physics Engine ASIC, title + 9 bullets) is ~136pt, so 140pt
% covers every entry. An absolute length is required because \baselineskip is 0
% in vertical mode, where \entrysep runs.
%
% A penalty-based alternative — wrapping each entry so no breakpoint exists
% inside it — was tried and rejected: with no legal breakpoint TeX produced an
% overfull \vbox and silently pushed ~30 lines off the last page. Never trade a
% split for lost content.
%
% Budget: this can leave up to ~#1 free at a page bottom, so the CV runs 4 pages
% rather than 3 at the master density. Accepted — this CV is not length-bound,
% and a split entry is the worse defect.
\newlength{\cvkeepskip}
\newcommand{\cvkeep}[1]{%
  \begingroup
    \setlength{\cvkeepskip}{#1}%
    \ifdim\pagegoal<\maxdimen
      \addtolength{\cvkeepskip}{\pagetotal}%
      \ifdim\cvkeepskip>\pagegoal
        \newpage
      \fi
    \fi
  \endgroup
}

% Separator between CV entries: uniform gap plus the keep-together guard.
\newcommand{\entrysep}{\vspace{2pt}\cvkeep{140pt}}

% --- Section heading spacing -------------------------------------------------
% \section*{} is used throughout: no "1 2 3" numbering, which a CV does not
% want, and pandoc renders it as an unnumbered <h2>.
%
% \cvsection wraps it so the surrounding vertical space can be tightened in one
% place. article.cls gives \section* roughly 3.5ex above and 2.3ex below, which
% across five sections costs about 12 lines — measured as 41 content lines on
% page 1 versus 48 for the master CV this mirrors.
%
% The \vspace{} calls live INSIDE a macro, which matters for two reasons:
%
%   1. Pandoc-safe. Tuning the same space via \makeatletter / \@startsection is
%      the obvious approach and silently BREAKS THE HTML VIEW: pandoc's LaTeX
%      reader then fails to recognise every \section* in the file and emits no
%      headings at all. Verified in isolation — identical \section*{} lines
%      inside a \makeatletter block produce zero <h1>/<h2>, while the same lines
%      inside this macro produce them correctly. cv/build.sh's "<h2" assertion
%      catches that failure, which is precisely why the check exists.
%   2. A macro keeps the spacing invisible to pandoc as content, so it cannot
%      leak into the HTML as stray spacing text or an empty paragraph.
%
% A negative \vspace is used rather than a package: titlesec/titletoc live in
% texlive-latex-extra, which this file deliberately avoids (see the
% dependency note in the preamble).
%
% `\vspace*` (starred) rather than plain `\vspace` before the heading: LaTeX
% discards a negative \vspace that is preceded by a \section's own
% \addvspace scaffolding, so the plain form silently did nothing. The starred
% form applies unconditionally. Measured: 26pt of space around each heading
% with the plain form, versus the intended tighter spacing with the starred one.
%
% The values below are the least aggressive that still look deliberate rather
% than cramped. Do not tighten further without re-checking legibility — the aim
% is a readable CV, not a short one.
\newcommand{\cvsection}[1]{%
  \vspace*{-11pt}%
  \section*{#1}%
  \vspace*{-10pt}%
}

% --- Date / meta span ------------------------------------------------------
% Wrapper for the dates and other trailing metadata on an entry.
%
% Why not just \emph{} or bare text: this file renders to HTML as well as PDF,
% and pandoc DROPS \hfill entirely, so `\textbf{Role} \hfill 2027` becomes
% `<strong>Role</strong> 2027` in the on-page view — the date loses its
% separation from the title with nothing left to style. \textsf{} survives the
% conversion as `<span class="sans-serif">`, which gives the HTML a hook
% (`.cv-doc .sans-serif` in assets/css/main.scss) and keeps the PDF unchanged.
%
% Write dates as:  \textbf{ROLE} \hfill \cvwhen{START - END}, using a plain
% ASCII hyphen for the range. Do NOT use LaTeX's "--": it typesets an en dash
% (U+2013), which is no longer plain ASCII in the text layer and can defeat
% literal date matching in an ATS. Verified: "Aug 2026 -- Present" extracted as
% U+2013, while "Aug 2026 - Present" extracts as a hyphen.
\newcommand{\cvwhen}[1]{\textsf{#1}}%
% ---------------------------------------------------------------------------

\pagestyle{empty}

\begin{document}
\pdfinterwordspaceon

% ---- Header ---------------------------------------------------------------
% Contact details are written as PLAIN TEXT, separated by \textbullet, rather
% than as \href mailto:/tel: links. Two reasons:
%
%   1. ATS parsers match on literal text. A hyperlink can be extracted as
%      "email me" while the address itself is lost, so the address is what must
%      appear in the text layer.
%   2. It mirrors the master CV's own header, so the two documents agree.
%
% The website and GitHub ARE linked, since a URL is both a link and literal text.
\begin{center}
  {\Large\bfseries Alden Grover}\\[2pt]
  Pittsburgh, PA \textbullet\ aldengrover@gmail.com \textbullet\ 646-539-8055\\
  \href{https://alden-g878.github.io}{alden-g878.github.io}
  \textbullet\
  \href{https://github.com/alden-g878}{github.com/alden-g878}
\end{center}

% ---- Summary --------------------------------------------------------------
\cvsection{Summary}
Seeking spring/summer 2027 internships in digital circuit and RTL design,
verification, and research. BS/MS combined student with Ph.D. intent, open to
industry roles between degrees.

% ---- Education ------------------------------------------------------------
\cvsection{Education}
% Institution and location are deliberately on ONE line separated by a comma,
% not split with \hfill. The master CV formats it this way, and it keeps the
% string "Carnegie Mellon University, Pittsburgh, PA" intact in the text layer
% for an ATS to match — \hfill would separate the two with a wide gap.
\textbf{Carnegie Mellon University, Pittsburgh, PA} \\
BS in Electrical and Computer Engineering, expected graduation December 2026 \\
Expected MS in Electrical and Computer Engineering (combined MS/BS program),
graduating 2028, with Ph.D. intent \\
\textbf{Honors:} College of Engineering Dean's List (Spring 2024) \\
\textbf{Notable Classwork:} 18-447: Introduction to Computer Architecture, 18-320: Microelectronic Circuits, 18-340: Hardware Arithmetic
    for Machine Learning, 18-349: Introduction to Embedded Systems

% ---- Technical Skills ------------------------------------------------------
\cvsection{Technical Skills}
\begin{itemize}\tightlist
  \item \textbf{RTL Design:} SystemVerilog, Verilog, ASIC RTL design and tape
    out, FPGA design, TinyTapeout, network-on-chip (NoC) design, USB 2.0
    hardware, finite state machines (FSMs) for datapath control, VGA waveform
    generation
  \item \textbf{RTL Verification and Simulation:} SystemVerilog testbenches,
    Cocotb, RTL waveform debug, formal reachability proofs for FSMs, static
    timing analysis (STA), constrained random testing
  \item \textbf{EDA and Flow:} Cadence Virtuoso, Cadence synthesis and
    simulation, Yosys, SPICE, RTL-to-GDSII flow (netlist, place and route,
    GDSII), TSMC 28nm planar PDK, Sky130 PDK, generic 45nm planar PDK, KiCad
  \item \textbf{Computer Architecture:} RISC-V (RV32I/M) ISA and microcode,
    pipelining, branch prediction, cache management and bus snooping
  \item \textbf{Clocking and Custom Circuits:} PLLs, ring oscillators, digitally
    controlled delay lines (DCDL), standard-cell library design, PPA (power,
    performance, area) optimization
  \item \textbf{Embedded Systems:} Embedded C++, RTOS development, bootloader
    design, UART/USART and QSPI, custom serial protocols, hard real-time
    systems, memory-mapped peripheral firmware development, sensor fusion
    (Kalman filter)
  \item \textbf{Hardware Acceleration:} Systolic arrays, neural-network
    inference, KNN acceleration, decision-tree acceleration, floating-point
    arithmetic
  \item \textbf{Languages:} C, C++, SystemVerilog, Python, Julia
  \item \textbf{Tools:} Git and GitHub, PCB design, CAD, CAM
\end{itemize}

% ---- Research Experience --------------------------------------------------
\cvsection{Research Experience}

\textbf{Design space exploration for Pareto-optimal design of a Digitally
Controlled Delay Line utilizing current starved and path choice methods} \\
Ken Mai and Jim Bain Research Group, Carnegie Mellon University \hfill
\cvwhen{Aug 2026 - Present (in preparation)}
\begin{itemize}\tightlist
  \item Leading investigation into a Pareto-optimal method to combine path
    choice and current starved Digitally Controlled Delay Line techniques.
  \item Investigating impact of existing implementations of path choice and
    current starved DCDL lines on range/resolution per unit area for PVT
    corners.
  \item Developing techniques for determining the optimal combination of path
    choice and current starved implementations.
  \item Verifying proposed techniques with brute force simulation on a
    small-scale model and directed simulation on a large-scale implementation.
\end{itemize}

\entrysep
\textbf{IC I/O Harness on Commercial Processes and Tools} \\
Ken Mai and Jim Bain Research Group, Carnegie Mellon University \hfill
\cvwhen{Nov 2025 - Present}
\begin{itemize}\tightlist
  \item Designed a portion of a taped-out ASIC with a focus on PLL, ring
    oscillators, and FPGA fabric in the TSMC 28nm PDK.
  \item Designed a ring oscillator and digital phase-locked loop that generates a
    320 MHz on-chip clock from a 20 MHz reference, a 16-fold increase.
  \item Redesigning the PLL for a new tape-out, targeting a 400 MHz on-chip
    clock frequency.
  \item Verifying PLL behavior with SPICE simulation.
  \item Designed custom FPGA fabric to interface with the I/O Harness, helping
    students familiarize themselves with it without the complicated pitfalls
    of ASIC design.
  \item Designed a PCB carrier board to interface the finalized ASIC with an
    external FPGA for bring-up and post-silicon validation.
  \item \textbf{Technologies:} RTL design in SystemVerilog, synthesis and
    simulation with Cadence.
\end{itemize}

\entrysep
\textbf{Nanopositioning System for Lithography and Van der Waals
Heterostructure Construction} \\
HackerFab, Carnegie Mellon University \hfill \cvwhen{Sep 2025 - Dec 2025}
\begin{itemize}\tightlist
  \item Led a team that improved high-accuracy nanopositioning system resolution
    from 10µm to 400nm.
  \item \textbf{Technologies:} Embedded C++, CAD and CAM, piezoelectric
    nanopositioning.
\end{itemize}

% ---- Project Experience ---------------------------------------------------
% ---- Leadership ------------------------------------------------------------
% Placed between Research and Project Experience on purpose: leadership is a
% supporting signal, so it must not interrupt the technical narrative, but it
% deserves a heading of its own rather than being buried in Technical Skills.
%
% Teaching is CMU coursework support (98-341), NOT a TA-ship for a different
% institution, and the two roles are the same course at different levels of
% responsibility — teacher now, TA previously. Listed as two dated entries
% rather than one, so the progression from TA to instructor is visible.
\cvsection{Leadership}

\textbf{Class Teacher, 98-341: Build Your Own Breadboard Computer} \\
Carnegie Mellon University \hfill \cvwhen{Aug 2026 - Present}
\begin{itemize}\tightlist
  \item Teaching 98-341, covering the design and construction of a
    breadboard computer.
\end{itemize}

\entrysep
\textbf{Teaching Assistant, 98-341: Build Your Own Breadboard Computer} \\
Carnegie Mellon University \hfill \cvwhen{Jan 2026 - May 2026}
\begin{itemize}\tightlist
  \item Supported students with hardware debugging and design questions in
    the same course.
\end{itemize}

\entrysep
\textbf{Electronics Master, Robotics Club} \hfill \cvwhen{Aug 2025 - Present}
\begin{itemize}\tightlist
  \item Lead the electronics subgroup of the club.
\end{itemize}

\cvsection{Project Experience}

\textbf{Standard Cell Library Design} \hfill \cvwhen{Apr 2026 - Present}
\begin{itemize}\tightlist
  \item Designed a standard cell library in a 45nm planar PDK.
  \item Extracted parasitic behavior for use in timing characterization.
  \item Verified functionality and characterized performance with SPICE
    simulation.
  \item Cell library consisted of NAND, NOR, and NOT gates as well as DFF and
    latch circuits.
\end{itemize}

\entrysep
\textbf{Physics Engine ASIC} \hfill \cvwhen{Dec 2025 - Present}
\begin{itemize}\tightlist
  \item Tape out: Sky130 open-source PDK (SkyWater 130nm), via the TinyTapeout shuttle.
  \item Implemented a falling sand simulation with output over VGA.
  \item QSPI memory controller reduced required on-chip memory to 1/8 of
    original size (from 96 bytes to 12 bytes).
  \item Complete ASIC design flow: RTL design, RTL simulation/debug, netlist
    generation, place and route, GDSII generation.
  \item Reduced clock cycles per physics simulation tick by a factor of three
    by adapting a physics simulation technique from graphics programming.
  \item Utilized RTL waveform debugging to fix bugs in the RTL design.
  \item Interfaced with external RAM hardware via QSPI, requiring intra-clock
    edge timing control.
  \item Leveraged fixed-delay digital buffers to ensure all protocols are
    obeyed.
  \item Constructed a test bench using Cocotb and Python, and verified
    functionality with assertions and constrained random testing.
\end{itemize}

\entrysep
\textbf{Lightweight Autonomous Vehicle Firmware, Robobuggy} \hfill
\cvwhen{Jan 2023 - Present}
\begin{itemize}\tightlist
  \item Designed, constructed, and maintained a high-performance embedded
    system in C++ over the course of four years. Lead designer for two of those
    years.
  \item Leveraged an existing embedded system architecture and carefully chosen
    consumer hardware to maximize performance per dollar.
  \item Embedded software engineering with hard real-time requirements and
    tight control loops. Average control loop time under 10 microseconds.
  \item Designed PCBs to interface microcontrollers with peripheral devices
    including GPS, servo actuators, accelerometers, encoders, and 433MHz radios.
  \item Collaborated with control-system experts to design an Unscented Kalman
    Filter to integrate disparate sensor streams into a cohesive whole.
\end{itemize}

\entrysep
\textbf{Static Timing Analysis Tool} \hfill \cvwhen{Jan 2026 - May 2026}
\begin{itemize}\tightlist
  \item Built a static timing analysis tool for the Yosys Open Synthesis Suite.
  \item Able to determine timing behavior of non-combinational loop netlists.
  \item Timing calculations using PVT corners ensure accuracy of results in all
    expected environmental conditions.
\end{itemize}

\entrysep
\textbf{7-stage Patterson and Hennessy style RV32I Core} \hfill
\cvwhen{Jan 2026 - May 2026}
\begin{itemize}\tightlist
  \item Implemented an RV32I core in SystemVerilog with 7 pipeline stages.
  \item Maximal forwarding between pipeline stages minimized data-dependent
    stalling. The final implementation was more than 5 times faster than the
    naive approach using real-world RV32I programs.
  \item Designed a two-bit branch predictor with a branch target buffer that
    doubled accurate branch predictions from 40\% to 80\%.
  \item Implemented a cache-management solution that successfully arbitrated
    cache conflicts inherent in non-atomic bus snooping systems.
\end{itemize}

\entrysep
\textbf{FP32 FPU and Neural Network Systolic Array} \hfill
\cvwhen{Sep 2025 - Dec 2025}
\begin{itemize}\tightlist
  \item Designed IEEE 754 compatible hardware for computing 32-bit floating
    point operations.
  \item Supports all arithmetic operations (add, subtract, multiply, divide).
  \item Performed mathematical modeling to optimize the addition system for
    power, performance, and area.
  \item Custom multiplication and division architecture enabled single-cycle
    operation and low-impact change to a highly pipelined design.
  \item Implemented the FP unit in a systolic array used to accelerate
    inference of a deep neural network.
\end{itemize}

\entrysep
\textbf{K Nearest Neighbor Inference Accelerator} \hfill
\cvwhen{Sep 2025 - Dec 2025}
\begin{itemize}\tightlist
  \item Designed a series of hardware threads that implemented the K Nearest
    Neighbor machine learning algorithm.
  \item Heavily pipelined design and Manhattan distance calculation maximized
    clock frequency.
  \item Exceeded the 300 MHz timing requirement (PDK specified by 18-340).
\end{itemize}

\entrysep
\textbf{Decision Tree Training Accelerator} \hfill \cvwhen{Sep 2025 - Dec 2025}
\begin{itemize}\tightlist
  \item Designed a series of hardware threads that accelerated decision tree
    machine learning algorithm training.
  \item Calculated Gini score using custom multiplication hardware.
  \item Leveraged arithmetic properties of numerical representation to bypass
    expensive computations.
  \item Custom multiplication hardware reduced cycles per operation from 32
    to 1.
\end{itemize}

\entrysep
\textbf{Network-on-a-Chip} \hfill \cvwhen{Jan 2025 - May 2025}
\begin{itemize}\tightlist
  \item Designed and simulated four port packet routing hardware supporting
    interconnected routers and fair packet flow between ports.
  \item Parallel FIFO buffers and internal counters ensure high throughput
    despite large data flow.
  \item Address logic enabled packet transmission through multiple routers to
    reach the destination.
  \item Internal data representation minimized packet latency through the
    router.
\end{itemize}

\entrysep
\textbf{USB 2.0 Hardware} \hfill \cvwhen{Jan 2025 - May 2025}
\begin{itemize}\tightlist
  \item Implemented read and write operations.
  \item 8-way pipelined and parallelized design enables rapid data transmission
    and reception, 8 times faster than a serial design.
  \item Custom orchestration hardware enables device-to-device transfers
    without a host OS.
\end{itemize}

\entrysep
\textbf{Matrix Multiplication, Addition, and Inner Product Hardware} \hfill
\cvwhen{Sep 2024 - Dec 2024}
\begin{itemize}\tightlist
  \item Interfaced with DSP slices and ROMs on an Intel Altera DE10 FPGA.
  \item Coordination between all available DSP hard cores and ROM modules
    within 98\% of the dataflow limit on the number of compute cycles.
\end{itemize}

\entrysep
\textbf{CPU Arithmetic Coprocessor and Microcode Modification} \hfill
\cvwhen{Sep 2024 - Dec 2024}
\begin{itemize}\tightlist
  \item Modified a multicycle CPU to support double-word arithmetic operations
    via microarchitecture changes to interface with an arithmetic coprocessor.
  \item Coprocessor design minimized instruction decode logic and microcode
    architecture, and eliminated software-based data dependence hazards using
    hardware solutions.
\end{itemize}

\entrysep
\textbf{Real-Time Operating System and Bootloader for ESP32} \hfill
\cvwhen{Jan 2024 - May 2024}
\begin{itemize}\tightlist
  \item Wrote bootloader and fixed priority preemptive scheduling RTOS for
    ESP32.
  \item Developed firmware for memory mapped and GPIO-based peripherals
    including USART/UART, H-Bridge BLDC motor controllers, rotational encoders,
    and LCD displays.
\end{itemize}

\end{document}
